\chapter{Conclusion and Discussion}
\label{ch:conclusion}

\section{Conclusion}
\label{sec:conclusion}

This report demonstrated that encoding tau-theory---the insect-inspired time-to-contact
regulation strategy of Lee \cite{lee1976}---as an observation and reward shaping signal
significantly improves the landing behaviour of a PPO-trained drone agent.

The tau agent achieves a \textbf{45\% reduction in touchdown speed}
($0.512 \rightarrow 0.281$~m/s) and an 8.9\% improvement in tau-regulation fidelity
compared to a baseline agent with no bio-inspired inductive bias.
Both improvements are statistically significant at $p < 0.001$ across 90 paired
evaluation episodes (Wilcoxon signed-rank test).
Crucially, these gains are fully preserved under zero-shot wind disturbances up to
2.0~m/s$^2$, demonstrating that the bio-inspired strategy is robust to domain shift.
A sensitivity analysis over six shaping coefficients identifies $\lambda = 0.10$ as
the optimal value, revealing a clear V-shaped trade-off between tau-regulation fidelity
and task success.

Analysis of the learned strategies shows that the baseline agent develops a
quasi-tau-regulation approach as an emergent property of optimising soft landing,
while the tau agent's explicit percept and shaping signal push it closer to the
biological optimum of $\dot{\tau} = -0.5$.
The persistent gap between the learned $\dot{\tau}$ distributions ($-1.59$ vs.\ $-1.75$)
and the biological target ($-0.5$) is explained by the competing time penalty, which
incentivises faster descent.

\section{Limitations}
\label{sec:limitations}

The following limitations should be noted when interpreting these results:
\begin{itemize}
  \item The environment is a 2-D point-mass simplification: attitude dynamics, rotor
        inertia, and 3-D wind are not modelled.
  \item The sensitivity sweep used $500\,000$ training steps (half the main budget),
        which may underestimate the performance of sub-optimal $\lambda$ values at full
        convergence.
  \item Only one seed was used for the sensitivity sweep; inter-seed variance at
        non-optimal coefficients may be larger than shown.
  \item The wind disturbance model (horizontal only) is a simplification;
        vertical gusts or turbulence would be a more realistic test.
\end{itemize}

\section{Future Work}
\label{sec:future}

Several directions are worth exploring in follow-up work:
\begin{itemize}
  \item Extending tau-regulation to a 3-D environment with full attitude dynamics,
        more closely matching real UAV flight.
  \item Using neuromorphic (event-based) vision sensors \cite{clady2014} that directly
        encode optic flow in the spiking domain, providing a more biologically faithful
        $\tau$ percept.
  \item Comparing PPO with evolutionary strategies \cite{such2017}, which may explore
        tau-regulation behaviour differently.
  \item Evaluating sim-to-real transfer using a physical quadrotor platform to assess
        whether the learned tau strategy holds under real aerodynamic effects.
\end{itemize}
